\cotefile{64}
\scanpage{91}
les \deleted{\unsure{objets|objet}} épinglages de Legendre-Jacobi ``directs'', i.e. tels que la section de $\mu_4^\times$ qu'\deleted{\unsure{elles}}il définit \unsure{soit|sont} justement $i$. \deleted{\unsure{leurs|Leurs|base}} \deleted{Les}Le \inserted{champ des} courbes elliptiques relatives \inserted{(sur des $S/\widetilde S_0$)} munies d'un épinglage de Legendre-Jacobi direct, est représentable par $U_{03\,\widetilde S_0}$, sur lequel on fait opérer le groupe modulaire $\mathbb D'_3$, qui est ainsi \inserted{justement} le sous-groupe de $\mathbb D_3\times\mu_{4\,\widetilde S_0}$ qui fait passer d'un épinglage direct à un épinglage direct.

Revenant au cas de la base absolue $\unsure{S_0|L_0}$, on doit encore expliciter les opérations de $\mu_4\times\unsure{\mathfrak S_3|\mathcal G_3}$ sur $\underline O_{U_{0,3}\times\mu_4^\times}\deleted{\unsure{(-1)|(1)}}(1)$, via l'action tautologique résultant de l'action de $\mathfrak S_3$ sur $V_{\{0,1,\infty\}}(\underline O_{S_0})$ d'une part, de \inserted{$\mathfrak S_3$ et} $\mu_4$ sur le facteur $\mu_4^\times$ d'autre part. On désigne l'action tautologique d'un $(g,\lambda)$ sur $\zeta\in\Gamma\underline O_{\unsure{U_{0,3}|\mathbb P_{0,3}|U_{03}}}(1)(\unsure{1|t|A})$
% archive sheet number: 90
